\section{The Order of Things}

\begin{quotex}
There are four nutriments for the maintenance of beings who have come into being or for the support of those in search of a place to be born: Physical food, gross or refined; touch as the second; thinking the third; and consciousness the fourth. \flright{\textsc{Buddha}, \emph{Samyatta Nikaya}}

\end{quotex}
\paragraph{Namaste}
In certain circles, ``Namaste" has replaced a hale ``Hello" as a form of greeting. Sometimes the English translation follows: ``The divine in me honors the divine in you." That may have some merit if taken literally, but the subtext is usually something like this: ``The vulgar part of me accepts the vulgar part of you."

What is less well known is that the major spiritual Tradition of the West has its own understanding of Namaste, i.e., the true understanding without the subtext:

\begin{quotex}
One loves one's neighbor for what there is of God in him and not for what he puts of himself in himself, that is, his sins, his whims, his personal ideas. We love him in the measure that he is with God. \flright{\textsc{Pope St. Pius X}}

\end{quotex}
\paragraph{Reality of Ideas}
Curiously, the writer \textbf{Charles Upton} grasped the reality of Ideas from the channeled Seth material. He quotes this passage:

\begin{quotex}
The term, a supreme being, is in itself distortive, for you naturally project the qualities of human nature upon it. If I told you that God was an idea, you would not understand what I meant, for you do not understand the dimensions in which an idea has its reality, or the energy it can originate and propel. You do not believe in ideas in the same way that you believe in physical objects, so if I tell you that God is an idea, you will misinterpret this to mean that God is less than real — nebulous, without reality, without purpose, and without motive action. 

\end{quotex}
From this, Upton realized that objective Ideas are not abstractions but rather densely-packed, conscious realities. He mentions that the Islamic theosophist Suhrawardi represents the Platonic Ideas as angels. Of course, this is also more or less the esoteric teaching of the West. The order of the angels corresponds to different levels of thought. Hence, our thinking is a participation in the transcendent order, good, bad, or neutral. The alternative is to believe that bio-electrochemical processes somehow create thoughts.

\paragraph{Psychic Reality in Jung}
Carl Jung likewise accepts ideas as entities to be studied by psychology, just a natural scientists study plants and animals. Although qua scientist he remains agnostic about the metaphysical reality of the ideas, that need not deter us. He explains:

\begin{quotex}
For modern psychology, ideas are entities, like animals and plants. The scientific method consists in the description of nature. All mythological ideas are essentially real, and far older than any philosophy. Like our knowledge of physical nature, they were originally perceptions and experiences. Insofar as such ideas are universal, they are symptoms or characteristics or normal exponents of psychic life, which are naturally present and need no proof of their ``truth". 

\end{quotex}
Universal ideas are part of the natural constitution of the psyche. Such ideas that are not consciously experienced are buried in the unconscious and may appear as neurotic symptoms. If you don't think that possible, recall that the Buddha distinguished thinking from consciousness.

For example, the idea of God is a psychic fact. The idea of immortality is a psychic fact, too. Hence, it is normal to think about it and abnormal not to.

Jung also researched ``so-called occult phenomena", although we prefer the more traditional expression, ``preternatural phenomena", including amnesia, dual consciousness, automatic writing, channeling, and so on. He pointed out that in practice it can be exceedingly difficult to distinguish such states from normal psychology and ``even with the psychology of the supranormal, that of genius." I mention this to explain how the Seth material can display signs of genius, as Upton noted, while being so wrong in other aspects.

\paragraph{Psychic Reality in Freud}
Freud is much maligned today but, like Jung (or Seth), he may have had some good ideas, provided we don't draw any metaphysical conclusions from his psychological observations. An oddity of the official English translation is that it replaces the concrete German pronouns with abstract Latin ones. Hence, instead of the ego and the id, we should rather understand ``the I" and ``the it". In that case, the I is never an object, but always a subject, and the It is an impersonal mass of impulses and desires. In terms of German idealism, the It is the not-I.

However, the contemporary mind understands things in quite the opposite way. The It has become the ``true self", since all such desires define who we really are. Then the I (or ego) becomes the object that acts on those desires.

In Freudian terms, the Pleasure Principle has achieved dominance over the Reality Principle. There is little point to appeal to reality and logic, when desire has become the ultimate arbiter of behavior. We see this in Freuds description of sexual perversion. By that term, Freud did not intend to imply a moral pejorative, but rather a neutral scientific term. ``Perversion" is simply the turning away from the biological function of sex. He wrote:

\begin{quotex}
The common characteristic of all perversions is that they have abandoned reproduction as their aim. We term sexual activity perverse when it has renounced the aim of reproduction and follows the pursuit of pleasure as an independent goal. 

\end{quotex}
Of course, that was the most common understanding up until a century or so ago. It was the teaching not only of the Catholic Church, but basically all the denominations. Freud conceded that touching, kissing, or licking were not perverted if they were part of foreplay prior to intercourse with the aim of reproduction. In our time, however, what was foreplay is now the end game, and reproduction is merely a risky side effect of sexual activity.

\paragraph{Old Age}
\begin{quotex}
Ordinarily we cling to our past and remain stuck in the illusion of youthfulness. Being old is highly unpopular. Nobody seems to consider that not being able to grow old is just as absurd as not being able to outgrow child-size shoes. A still infantile man of thirty is surely to be deplored, but a youthful septuagenarian—isn't that delightful? And yet both are perverse, lacking in style, psychological monstrosities. A young man who does not fight and conquer has missed the best part of his youth, and an old man who does not know how to listen to the secrets of the brooks, as they bumble down from the peaks to the valleys, makes no sense; he is spiritual mummy who is nothing but a rigid relic of the past. He stands apart from life, mechanically repeating himself to the last triviality. \flright{\textsc{Carl Jung}}

\end{quotex}
Yet we promote the spiritual mummies. We praise retirees who run a marathon, 60 year old supermodels who wear a bikini, or a grandmother foolishly playing virtual reality Star Wars games. There are many baby boomers who have not had an original thought since they were 19 years old.

\paragraph{Barnyard Animals}
\begin{quotex}
We must not forget that to speak of blood in the case of a man is not the same thing as to speak of it in the case of an animal. If, by blood, one means the biological heredity of a race, then in the animal, race is everything, while, in man, it is only a part. The error of certain race fanatics, who think that the reintegration of a race in its ethnic unity signifies \emph{ipso facto} the rebirth of a people, lies exactly here: they regard man as if he could be regarded as horses, cats, or dogs of a ``pure breed". The preservation or the restoration of the purity of race, in the narrowest sense, can be everything in an animal, but not in man—in the man of superior type: even for man, it can constitute a condition which may be necessary under certain aspects, yet is not sufficient in any case, since the racial factor is not the only one which defines man. 

\end{quotex}
That is Julius Evola's description of man and his relationship to biology. Somehow that is considered evil today. Nevertheless, the ``educated" scientific view is that man is indeed just a biological creature, not fundamentally different from horses, cats, and dogs. Go figure.

\paragraph{The Opposite of a Revolution}
There is a persistent misunderstanding of what the terms left and right might mean in their socio-political context. A careful reading of \emph{The Communist Manifesto} by \textbf{Karl Marx}, however, will bring much clarity. These are its two fundamental points; everything else follows upon them:

\begin{enumerate}
\item A spectre is haunting Europe — the spectre of communism. 
\item In short, the Communists everywhere support every revolutionary movement against the existing social and political order of things. 
\end{enumerate}
Those statements are from the beginning and the end of the Manifesto. The fundamental impulse of Marxism comes from a ``spectre", that is, an evil spirit. You may prefer the euphemism that it arises from an ``idea", or a meme. In that case, it is easy to forget that ideas are ``conscious realities"; a fortiori, it neglects to take Marx's own claim seriously.

The goal of Marxism, as Marx pointed out, is to support ``the forcible overthrow of all existing social conditions", i.e., the socio-political order (not just the economic order), wherever and whenever possible. This is always couched in terms of increased freedom or liberation, so that even many who claim to oppose Marxism, are actually supporting its aims.

Once this is grasped, contemporary events reveal themselves. There is no rationality to the revolution, since it is guided by evil spirits. The revolution always denies the Logos, and attacks Him constantly.

\textbf{Joseph de Maistre} famously claimed that what is required is the ``opposite of a revolution", that is, the establishment of the proper socio-political order of things. That can be understood as the recovery of the values that ``every well-bred man considered healthy and normal" prior to the revolution. This is not nostalgia for a return to the social and material conditions of a bygone era. The natural, or better, supernatural, order is transcendent to the vagaries of history. Thus, Tradition has manifested in many different concrete forms, while retaining its transcendent values. Marxism, on the other hand, is a materialist understanding of history.

In sum, every counter-revolutionary movement must have a deep understanding of the proper order of things; otherwise, it will be part of the problem, not part of the solution.

\paragraph{The Regression of Castes}
Rene Guenon has documented the regression of the castes: the spiritual authority gives way to political power, and so on. Curiously, Marxism and Tradition both agree on the transfer of power from caste to caste, or from class to class in Marxist terms. They only disagree on its significance.

There is a cartoonish diagram of Evola's explanation of the regression that has been making the rounds. Presumably, it is for those who cannot follow a written text, as it gets a lot of attention on social media. What would be more helpful is a little quiz that would identify which caste you belong to, based on your personality, character, abilities, and level of understanding. There will be a lot of suprises.

\paragraph{Planet of the Apes}
I found my copy of \emph{Fads and Fallacies in the Name of Science}, by \textbf{Martin Gardner}, which I had first read as a boy. The cranks he refutes are, unfortunately, more interesting than Gardner. A lot of chapters are on alternative medicine, the age of the earth, and anti-Darwinism.

One of the more interesting is Francis Crookshank\footnote{\url{https://en.wikipedia.org/wiki/Francis_Graham_Crookshank}}, who related human races to the apes. He claimed that the White race is like the chimpanzee, the Oriental like the orangutan, and the Black race like the gorilla. Now I've only seen the original \textit{Planet of the Apes}, and that was some time ago. Does Crookshank's theory offer a key to understanding the hidden meaning of the movie? When the apes someday film the movie, \emph{Planet of the Humans}, then we will know for sure.



\flrightit{Posted on 2017-12-29 by Cologero }

\begin{center}* * *\end{center}

\begin{footnotesize}\begin{sffamily}



\texttt{patriciakay3 on 2017-12-29 at 13:28 said: }

How helpful your writing has been over the years. You have brought about an order in my thinking that has centered my life in greater clarity. An inner ``lightedness" has emerged from a chaotic darkness, and you have contributed to this happening. Thank you.


\hfill

\texttt{james on 2017-12-30 at 04:30 said: }

Thanks for all the great posts over the years.Looking forward to whatever new format you adapt.


\hfill

\texttt{kalpatraveller on 2017-12-30 at 15:54 said: }

I am well aware you have a more important set of topics for the remainder of Gornahoor's incarnation (looking forward to all of them), than the one I am about to suggest, but it would be appreciated (at least by someone at my relatively unintegrated level) to know what kinds of things you might have found of value in the Seth material (the ``signs of genius"), as well as what was ``so wrong" in his case otherwise, perhaps channeling in general.

Rudolf Steiner once said, near the end of his life, that the time would come when many more would connect with higher realms, and while channeling might seem a likely example of that, I've had unresolved doubts. Richard Leviton, who was (and still seems to be) an indebted follower of Steiner, generally denounced the channeling phenomenon in his earlier book surveying anthroposophy, The Imagination of Pentecost. However, it seems more than ironic that he now practices a kind of ``clairvoyance" in which he communes with entities such as Merlin, and has regular contact with the Ofanim, his angelic ``mentors" at the level of the thrones…


\hfill

\texttt{Auld Wat on 2017-12-31 at 06:10 said: }

Thank you for another great post, Cologero!


\hfill

\texttt{Cologero on 2017-12-31 at 11:32 said: }

The Seth material has already been analyzed by Charles Upton in ``The System of the Antichrist", so I don't need to repeat it all here.

As for Steiner, start with this: \textit{Rudolf Steiner Gave Two Paths: The Science Path And The Occult Path:} \url{https://www.youtube.com/watch?v=ApMCQn0HXuM}

I am more interested in the philosophical (scientific) Steiner rather than the later clairvoyant (occult) Steiner.

For a analysis of clairvoyance vs creative imagination, check this out: \textit{The Theosophical Imagination:} \url{http://correspondencesjournal.com/wp-content/uploads/2017/12/16401_20537158_hanegraaff.pdf}


\hfill

\texttt{Malopello on 2018-01-03 at 05:30 said: }

Thanks for your blog, I discovered it recently, it's a pleasure to read.

I don't know if you already wrote about modern notions as Progress, or the mechanization of human life to a subhuman life, transhumanism, the drying of human imaginary, scientism, mass communication and information, etc.. And relating it to philosophers, saints, or whatever you want.

Some interesting readings that you may know :

The Reign of Quantity \& the Signs of the Times – by René Guénon

You seem to speak fluently french so this book too can be good too (didnt find in english) : la France contre les robots – by Georges Bernanos

I think it would be interesting to read your thoughts about it.

I dont know if this type of subjects can have its place in your blog, whatever thanks for your writings.


\end{sffamily}\end{footnotesize}
